% =====================================================================
%  FASE DE INICIALIZACIÓN (versión actualizada). Cambios respecto al original:
%   - Cloud Functions y App Check en la configuración, la arquitectura y la
%     estructura del proyecto; la clave de USDA ya no viaja en la aplicación.
%   - Reglas de seguridad vigentes (rutinas solo mediante función).
%   - Nueva subsección de consideraciones éticas y de privacidad.
%   - Flujo de ramas con las solicitudes de integración #5, #6 y #7.
% =====================================================================
% =====================================================================
\section{Fase de Inicialización}
\label{sec:des-inicializacion}

En esta fase se preparó el entorno técnico, se definió la arquitectura del
sistema y se diseñó el modelo de datos.

Durante la fase de Inicialización se establecieron las herramientas para implementar
la aplicación, gestionar el acceso y la persistencia de datos, e integrar
los componentes de reconocimiento de máquinas y generación de rutinas.

\subsection{Entorno de trabajo}

\begin{itemize}
    \item \textbf{React Native:} Framework para crear aplicaciones móviles con React.
    \item \textbf{TypeScript:} Lenguaje basado en JavaScript que incorpora tipado estático.
    \item \textbf{Expo:} Plataforma con herramientas para desarrollar, probar y compilar aplicaciones React Native.
    \item \textbf{Expo Router:} Sistema de navegación que define las rutas a partir de los archivos del proyecto.
    \item \textbf{Firebase:} Plataforma que ofrece servicios de autenticación, bases de datos, almacenamiento de archivos y funciones en la nube.
    \item \textbf{TensorFlow Lite:} Biblioteca para ejecutar modelos de aprendizaje automático en dispositivos móviles y embebidos.
    \item \textbf{Ultralytics YOLO:} Familia de modelos de visión artificial que permite detectar y localizar objetos en imágenes.
    \item \textbf{ARCore:} Plataforma de realidad aumentada que permite seguir el movimiento del dispositivo y anclar elementos virtuales al entorno físico.
    \item \textbf{USDA FoodData Central:} Base de datos sobre la composición nutricional de los alimentos.
    \item \textbf{Git:} Sistema de control de versiones que registra los cambios en los archivos.
    \item \textbf{GitHub:} Plataforma para alojar repositorios Git y colaborar en su desarrollo.
\end{itemize}

\subsection{Configuración del entorno de desarrollo}

En esta etapa se configuraron los recursos físicos y técnicos necesarios 
para preparar el entorno de desarrollo de aplicaciones para Android.

\subsection{Configuración para aplicaciones móviles android}

\begin{itemize}
    \item \textbf{Tipo de proyecto.} Aplicación móvil para Android desarrollada
    con React Native y TypeScript, sobre la plataforma Expo.

    \item \textbf{Librerías principales.}
    \begin{itemize}
        \item \archivo{expo-router} para organizar la navegación entre pantallas
        mediante rutas basadas en archivos.
        \item \archivo{firebase} para gestionar la autenticación, consultar y
        guardar datos en Cloud Firestore, acceder a Cloud Storage e invocar
        funciones del servidor.
        \item \archivo{expo-camera} y \archivo{expo-image} para capturar imágenes
        con el escáner y mostrar imágenes y GIF del catálogo.
        \item \archivo{react-native-fast-tflite} para ejecutar en el dispositivo
        los modelos de clasificación de máquinas y recomendación de rutinas.
        \item \archivo{react-hook-form} y \archivo{zod} para gestionar los
        formularios y validar los datos ingresados.
        \item \archivo{zustand} para gestionar el estado compartido de la
        aplicación y conservarlo entre sesiones.
        \item \archivo{react-native-chart-kit} para generar los gráficos del
        reporte de progreso físico.
    \end{itemize}

    \item \textbf{Configuraciones técnicas.}
    \begin{itemize}
        \item Generación de APK mediante EAS Build para distribución en Android[cite: 2].
        \item Activación de permisos de cámara.
        \item Configuración de reglas de seguridad para control de acceso a Cloud Firestore y Cloud Storage[cite: 2, 7].
        \item Definición de perfiles de compilación (\emph{development}, \emph{preview} y \emph{production}) en \archivo{eas.json} para automatizar el empaquetado en la nube[cite: 3].
        \item Uso opcional de funciones de servidor para procesar validaciones y tareas en la nube mediante Firebase[cite: 7].
    \end{itemize}
\end{itemize}


\subsection{Planificación inicial}
Durante esta etapa se establecerán los requisitos previos técnicos y funcionales de cada
proceso, teniendo en cuenta las necesidades del usuario y las condiciones de la infraestructura tecnológica.

\subsubsection{Análisis de procesos}

La implementación de cada proceso requiere que se cumplan las condiciones 
previas que se detallan a continuación:

{\small
\begin{longtable}{|L{1.3cm}|L{3.3cm}|L{9cm}|}
\hline
\rowcolor{encabezado}
\textbf{Código} & \textbf{Proceso} & \textbf{Pre-requisitos} \\
\hline
\endfirsthead

\hline
\rowcolor{encabezado}
\textbf{Código} & \textbf{Proceso} & \textbf{Pre-requisitos} \\
\hline
\endhead

P01 & Autenticación y control de acceso &
\begin{itemize}
    \item Configuración de Firebase Authentication para el acceso con correo
    electrónico y contraseña o con una cuenta de Google.
    \item Validación de las credenciales mediante el proveedor seleccionado.
    \item Creación del documento de perfil en Cloud Firestore durante el primer
    acceso.
    \item Verificación del estado de la sesión y del perfil para determinar
    qué pantallas puede abrir el usuario.
\end{itemize}
\\ \hline

P02 & Configuración inicial del perfil &
\begin{itemize}
    \item Usuario autenticado mediante el proceso P01.
    \item Formulario para registrar el objetivo, el nivel de entrenamiento,
    la edad, la estatura, el peso y los días disponibles.
    \item Validación de los campos obligatorios y de los valores ingresados.
    \item Guardado del perfil en Cloud Firestore antes de habilitar el acceso
    a las pantallas principales.
\end{itemize}
\\ \hline

P03 & Consulta de la guía de máquinas &
\begin{itemize}
    \item Catálogo de máquinas y ejercicios disponible para consulta.
    \item Datos descriptivos e imágenes asociados a cada máquina.
    \item Campo de búsqueda y filtros para localizar máquinas en el catálogo.
    \item Ficha con instrucciones de uso para la máquina seleccionada.
\end{itemize}
\\ \hline

P04 & Escaneo e identificación de máquinas &
\begin{itemize}
    \item Dispositivo Android con cámara y permiso de acceso concedido.
    \item Imagen de la máquina obtenida mediante el escáner.
    \item Modelo local de TensorFlow Lite disponible para clasificar la imagen.
    \item Conexión a Internet y acceso al servicio de detección cuando se
    utiliza el reconocimiento en la nube.
    \item Ficha disponible en el catálogo y opción de selección manual.
    \item Dispositivo compatible con ARCore para anclar la ficha al entorno
    físico cuando se utiliza esa modalidad.
\end{itemize}
\\ \hline

P05 & Generación de rutinas &
\begin{itemize}
    \item Usuario autenticado con el perfil completado en P02.
    \item Objetivo, nivel de entrenamiento y días disponibles registrados
    en el perfil.
    \item Modelo local de recomendación y catálogo de ejercicios disponibles
    para conformar la rutina.
    \item Acceso a la función \archivo{guardarRutina} para validar la rutina
    en el servidor y guardarla como activa.
\end{itemize}
\\ \hline

P06 & Planificación de rutinas &
\begin{itemize}
    \item Usuario autenticado con acceso al planificador.
    \item Selección de días de entrenamiento y ejercicios del catálogo.
    \item Registro de series, repeticiones y tiempos de descanso por ejercicio.
    \item Validación de la planificación antes de guardar la rutina
    personalizada mediante la función \archivo{guardarRutina}.
\end{itemize}
\\ \hline

P07 & Registro de alimentos &
\begin{itemize}
    \item Usuario autenticado y fecha de registro seleccionada.
    \item Texto de búsqueda para consultar alimentos mediante la función
    \archivo{buscarAlimentos}, conectada a USDA FoodData Central.
    \item Selección de un alimento con información nutricional disponible.
    \item Cantidad en gramos para calcular las calorías y los macronutrientes
    de la porción.
    \item Guardado de la porción en el diario de alimentación del usuario.
\end{itemize}
\\ \hline

P08 & Registro corporal &
\begin{itemize}
    \item Usuario autenticado con acceso al módulo de progreso físico.
    \item Peso en kilogramos y estatura en centímetros con valores válidos.
    \item Cálculo del índice de masa corporal y determinación de la categoría
    correspondiente.
    \item Guardado del registro con su fecha para incorporarlo al historial.
\end{itemize}
\\ \hline

P09 & Seguimiento de entrenamientos &
\begin{itemize}
    \item Usuario autenticado con una rutina activa.
    \item Selección de una sesión de entrenamiento planificada.
    \item Confirmación de que la sesión fue completada.
    \item Registro de la sesión y su fecha para actualizar la racha,
    el conteo semanal y el total de entrenamientos.
\end{itemize}
\\ \hline

P10 & Generación de reportes de progreso &
\begin{itemize}
    \item Usuario autenticado con registros corporales almacenados en P08.
    \item Consulta del historial en orden cronológico.
    \item Datos disponibles para mostrar la categoría actual y, cuando
    existen registros comparables, calcular la variación corporal.
    \item Componente de gráficos para representar la evolución del usuario.
    \item Componente de generación de PDF para exportar el reporte.
\end{itemize}
\\ \hline
P11 & Acceso a la aplicación &
\begin{itemize}
    \item Archivo .apk generado.
    \item Dispositivo compatible.
    \item Permisos de instalación de fuentes desconocidas.
    \item Espacio de almacenamiento suficiente en el dispositivo.
    \item Servidor \emph{backend} disponible 24/7 (Firebase Cloud).
\end{itemize}
\\ \hline
\caption{Cuadro de análisis de procesos (Elaboración propia)}
\label{tab:procesos}\\
\end{longtable}
}

\subsubsection{Planificación de las fases}

La funcionalidad se distribuyó en seis iteraciones de producción
(Tabla~\ref{tab:plan-iteraciones}), ordenadas de modo que cada una se apoye en
las anteriores: primero la identidad del usuario, luego el catálogo, que
alimenta al escáner y a las rutinas, y finalmente los módulos de seguimiento.

\subsubsection{Planificación de las fases}

El desarrollo se organizó según las fases de Mobile-D. La
Tabla~\ref{tab:plan-iteraciones} resume las actividades de cada fase y
la secuencia de las seis iteraciones de producción.

\begin{table}[H]
\centering
\small
\renewcommand{\arraystretch}{1.2}


\begin{tabularx}{\textwidth}{|L{2.8cm}|C{2.2cm}|Y|}
\hline
\rowcolor{encabezado}
\textbf{Fase} & \textbf{Iteración} & \textbf{Descripción} \\
\hline

\textbf{Exploración} & Iteración 1 &
Identificación de los interesados y actores del sistema, definición de
los requerimientos funcionales y no funcionales, y delimitación del alcance. \\
\hline

\textbf{Inicialización} & Iteración 2 &
Configuración del entorno de desarrollo para Android, análisis de los
procesos, planificación de las iteraciones y diseño de la arquitectura,
el modelo de datos y la interfaz de usuario. \\
\hline

\textbf{Producción} & Iteración 3 &
Implementación de la autenticación y el perfil de usuario, la guía de
máquinas y ejercicios, el escáner para el reconocimiento de máquinas,
la generación y planificación de rutinas, el diario de alimentación
y el seguimiento del progreso físico con reportes en PDF. \\
\hline

\textbf{Estabilización} & Iteración 4 &
Integración de los módulos, refactorización del código y verificación
del funcionamiento conjunto de la aplicación. \\
\hline

\textbf{Finalización} & Iteración 5 &
Preparación de la documentación final, el manual de usuario, los
informes de pruebas y la versión de la aplicación para su distribución
en formato APK. \\
\hline

\end{tabularx}
\caption{Cuadro de análisis de planificación de fases}
\label{tab:plan-iteraciones}
\fuente{Elaboración propia.}
\end{table}

\subsection{Diseño de la arquitectura del sistema}
\label{sec:des-arquitectura}

La arquitectura del proyecto sigue un modelo cliente--servidor. 
Como se muestra en la Figura~\ref{fig:arquitectura}, la aplicación móvil 
se comunica directamente con los servicios de Firebase para gestionar 
la autenticación, la base de datos y el almacenamiento, y consulta 
la API pública de wger (wger Project, s.~f.) para obtener información sobre ejercicios. 
Para consultar los datos nutricionales de USDA y reconocer máquinas 
mediante Ultralytics, la aplicación utiliza las funciones correspondientes 
de Cloud Functions.

\begin{figure}[H]
\centering
\begin{tikzpicture}[font=\small,
  lay/.style={rectangle, draw=black!60, fill=#1, minimum width=8.6cm,
    minimum height=1.25cm, align=left, text width=8.3cm, font=\footnotesize},
  srv/.style={rectangle, rounded corners=3pt, draw=black!60, fill=orange!8,
    minimum width=2.9cm, minimum height=0.95cm, align=center, font=\footnotesize}]
  % Cliente
  \node[lay=green!12] (vista) at (0,4.3) {\textbf{Vista} --- \archivo{app/}, \archivo{components/}\\
    Pantallas de Expo Router, componentes de interfaz, NativeWind};
  \node[lay=green!7] (ctrl) at (0,2.65) {\textbf{Controlador} --- \archivo{controllers/}, \archivo{utils/}, \archivo{store/}\\
    Casos de uso, reglas de negocio, generador de rutinas, estado global};
  \node[lay=green!4] (modelo) at (0,1.0) {\textbf{Modelo} --- \archivo{models/}, \archivo{schemas/}\\
    Entidades, esquemas Zod, repositorios (únicos que conocen Firestore)};
  \node[lay=black!4] (serv) at (0,-0.65) {\textbf{Servicios} --- \archivo{services/}\\
    Firestore, Storage, llamadas a Cloud Functions, cliente HTTP de wger, modelos de IA y módulo nativo de ARCore};
  \node[draw, dashed, rounded corners, fit=(vista)(serv), inner sep=8pt,
    label={[font=\small\bfseries]above:Cliente: aplicación Android}] (cli) {};
  \draw[flecha] (vista) -- (ctrl);
  \draw[flecha] (ctrl) -- (modelo);
  \draw[flecha] (modelo) -- (serv);
  % Servidor
  \node[srv] (auth) at (9.2,4.3) {Firebase\\Authentication};
  \node[srv] (fs) at (9.2,2.65) {Cloud\\Firestore};
  \node[srv] (st) at (9.2,1.0) {Cloud\\Storage};
  \node[srv] (cf) at (9.2,-0.2) {Cloud Functions\\+ App Check};
  \node[srv, fill=blue!6] (wg) at (9.2,-1.5) {APIs externas\\(wger, USDA)};
  \node[srv, fill=blue!6] (ul) at (9.2,-2.8) {Ultralytics\\(\emph{endpoint} YOLO)};
  \node[draw, dashed, rounded corners, fit=(auth)(ul), inner sep=8pt,
    label={[font=\small\bfseries]above:Servidor}] (sv) {};
  \draw[flecha, <->] (cli.east |- auth) -- node[above, font=\tiny]{SDK} (auth);
  \draw[flecha, <->] (cli.east |- fs) -- node[above, font=\tiny]{tiempo real} (fs);
  \draw[flecha, <->] (cli.east |- st) -- (st);
  \draw[flecha, <->] (cli.east |- cf) -- node[above, font=\tiny]{HTTPS} (cf);
  \draw[flecha, <->] (cf) -- node[right, font=\tiny]{USDA} (wg);
  \draw[flecha, <->] (cf.east) -- ++(0.5,0) |- node[right, pos=0.2, font=\tiny]{imagen} (ul.east);
  \draw[flecha, <->] (cli.east |- wg) -- node[above, font=\tiny]{wger} (wg);
\end{tikzpicture}
\caption{Arquitectura cliente--servidor y capas del cliente}
\label{fig:arquitectura}
\fuente{Elaboración propia.}
\end{figure}

\begin{table}[H]
\centering
\small
\begin{tabularx}{\textwidth}{|L{3.4cm}|L{2.6cm}|Y|}
\hline
\rowcolor{encabezado} \textbf{Capa conceptual} & \textbf{Capa MVC} & \textbf{Implementación} \\
\hline
Presentación & Vista & Pantallas y componentes de \archivo{app/} y
\archivo{components/}. \\ \hline
Lógica de negocio & Controlador & Controladores de \archivo{controllers/},
reglas puras de \archivo{utils/} (IMC, generador de rutinas, estadísticas de
entrenamiento). \\ \hline
Servicios (Firebase y APIs) & Modelo y Servicios & Repositorios de
\archivo{models/repositories/}; \archivo{firestoreService},
\archivo{storageService}, \archivo{functionsService} (llamadas a las funciones
en la nube), \archivo{wgerService} y \archivo{usdaService} (que invoca la función
\archivo{buscarAlimentos}). \\ \hline
Inteligencia artificial & Controlador y Servicios & Generador de rutinas
basado en reglas (\archivo{utils/}) y, en \archivo{services/ai/}, el
recomendador de rutinas y el reconocedor de máquinas, ambos modelos de
TensorFlow Lite ejecutados en el dispositivo, más la consulta opcional al
detector de la nube mediante \archivo{reconocerMaquina}. \\ \hline
Realidad aumentada & Vista y módulo nativo & Vista \archivo{CoachAR}
(\archivo{modules/coach-ar/}, Kotlin) y su interfaz tipada en
\archivo{services/ar/}; el controlador solo entrega a la vista la caja y la
ficha de la máquina confirmada. \\
\hline
\end{tabularx}
\caption{Correspondencia entre las capas conceptuales y las capas del cliente}
\label{tab:capas-mvc}
\end{table}

\subsubsection{Modelo de datos NoSQL}
\label{subsec:modelo-datos}

El modelo de datos NoSQL define cómo se organiza, almacena y consulta la 
información en una base de datos no relacional. Cloud Firestore utiliza 
un modelo documental en el que los datos se agrupan en colecciones. 
Cada colección contiene documentos, y cada documento almacena la información 
en campos.
A continuación el esquema diseñado para la base de datos del sistema.

\begin{figure}[H]
\centering
\begin{tikzpicture}[
  font=\footnotesize,
  documento/.style={
    entidad,
    text width=0.38\linewidth,
    align=left,
    inner sep=5pt,
    outer sep=0pt
  },
  grupo/.style={font=\footnotesize\bfseries},
  enlace/.style={semithick, draw=black!65}
]

% Colección de usuarios
\node[documento] (usr) {\texttt{usuarios/\{uid\}}
  \nodepart{second}
  nombre, apellidos, email\\
  genero, altura, pesoActual, edad\\
  diasDisponibles, nivel, objetivo\\
  rutinaActivaId\\
  creadoEn, actualizadoEn
};
\node[grupo, above=3pt of usr] {Perfil del usuario};

% Subcolecciones privadas, organizadas en dos filas
\coordinate[below=1.1cm of usr] (privadas);
\node[grupo] at (privadas) {Subcolecciones privadas};

\node[documento, anchor=north east] (rut)
  at ([xshift=-0.8cm,yshift=-0.6cm]privadas)
  {\texttt{rutinas/\{id\}}
    \nodepart{second}
    nombre, tipo, objetivo, nivel\\
    diasPorSemana, planSemanal[\,]\\
    origenPlan, versionModelo\\
    confianzaModelo, usuarioId, creadoEn
  };

\node[documento, anchor=north west] (reg)
  at ([xshift=0.8cm,yshift=-0.6cm]privadas)
  {\texttt{registrosNutricionales/\{id\}}
    \nodepart{second}
    \textbf{IMC:} weight, height, bmi, category\\
    \textbf{Ingesta:} alimento, calorias, cantidad\\
    proteinas, carbohidratos, grasas\\
    fdcId, usuarioId, creadoEn
  };

\coordinate (segundaFila) at ([yshift=-0.55cm]reg.south);
\node[documento, anchor=north east] (pro)
  at (rut.east |- segundaFila)
  {\texttt{progresoFisico/\{id\}}
    \nodepart{second}
    peso, altura, imc, notas\\
    usuarioId, creadoEn
  };

\node[documento, anchor=north west] (ent)
  at (reg.west |- segundaFila)
  {\texttt{entrenamientos/\{fecha\}}
    \nodepart{second}
    rutinaId, dia (1--7)\\
    fecha (= id), usuarioId, creadoEn
  };

% Eje central que conecta el perfil con sus cuatro subcolecciones
\draw[enlace] (usr.south) -- ([yshift=3pt]privadas.north);
\draw[enlace] ([yshift=-3pt]privadas.south) -- (privadas |- ent.west);
\draw[enlace] (rut.east) -- (privadas |- rut.east);
\draw[enlace] (reg.west) -- (privadas |- reg.west);
\draw[enlace] (pro.east) -- (privadas |- pro.east);
\draw[enlace] (ent.west) -- (privadas |- ent.west);

% Colecciones raíz compartidas
\coordinate (catalogos) at ([yshift=-0.75cm]privadas |- ent.south);
\node[grupo] at (catalogos) {Catálogos de solo lectura};

\node[documento, anchor=north east] (eje)
  at ([xshift=-0.8cm,yshift=-0.35cm]catalogos)
  {\texttt{ejercicios/\{id\}}
    \nodepart{second}
    nombre, grupoMuscular[\,]\\
    descripcion\\
    seriesRecomendadas\\
    repeticionesRecomendadas
  };

\node[documento, anchor=north west] (maq)
  at ([xshift=0.8cm,yshift=-0.35cm]catalogos)
  {\texttt{maquinas/\{id\}}
    \nodepart{second}
    nombre, categoria, tipo, descripcion, orden\\
    musculosPrincipales[\,], musculosSecundarios[\,]\\
    instrucciones[\,], erroresComunes[\,]\\
    prevencionLesiones[\,]\\
    imagenPrincipal, gifPrincipal\\
    variantes[\,], ejercicios[\,], ejerciciosGif[\,]
  };

\end{tikzpicture}
\caption{Modelo de datos en Cloud Firestore}
\label{fig:modelo-datos}
\fuente{Elaboración propia.}
\end{figure}

\subsubsection{Diseño de seguridad y controles}

El diseño de seguridad de la aplicación busca proteger los datos personales
y preservar la integridad de la información almacenada. Para ello, se
establecieron controles de autenticación, acceso a los datos y validación
de las operaciones que se ejecutan en el servidor.

\begin{itemize}
    \item \textbf{Base de datos.} Las reglas de seguridad de Cloud Firestore
    utilizan request.auth para identificar al usuario cuando se
    autentica, obtener el uid y la información solicitada.
    Estas reglas permiten definir quién puede leer, crear, actualizar o
    eliminar documentos. En la aplicación, se utilizan para:
    \begin{itemize}
        \item Restringir el acceso al perfil, las rutinas y los registros
        de alimentación, progreso físico y entrenamiento al usuario
        propietario de los datos.
        \item Permitir que los usuarios autenticados consulten los catálogos
        de máquinas y ejercicios, y reservar su modificación para
        operaciones autorizadas.
        \item Impedir que la aplicación cree rutinas directamente en
        Firestore. La función \archivo{guardarRutina} verifica la identidad
        del usuario y valida los datos antes de guardarlos.
    \end{itemize}

    \item \textbf{Almacenamiento de archivos.} Las reglas de Cloud Storage
    restringen la lectura de los recursos del catálogo a usuarios
    autenticados y controlan quién puede modificar los archivos.

    \item \textbf{Funciones del servidor.} Las funciones de búsqueda de
    alimentos, guardado de rutinas y reconocimiento de máquinas verifican
    la identidad del usuario. Firebase App Check, integrado con Play
    Integrity, complementa este control al verificar el origen de las
    solicitudes. Las claves de los servicios externos se mantienen en
    el servidor.
\end{itemize}

\subsubsection{Consideraciones éticas y de privacidad}
\label{sec:etica}

La aplicación almacena datos personales relacionados con la salud (peso,
estatura, edad, género y objetivo de entrenamiento) asociados a una cuenta de
correo electrónico. Asimismo, el sistema no reemplaza la supervisión de entrenadores ni de
nutricionistas, y las estimaciones de calorías, IMC y metas diarias son
referenciales. Para entrenar el reconocedor de máquinas se utilizaron imágenes
de catálogos y fotografías tomadas en distintos gimnasios del departamento
de Cochabamba.

Con el modo de reconocimiento de maquinas en la nube activado, la fotografía que toma el
escáner sale del teléfono: viaja a la Cloud Function y de ahí al \emph{endpoint}
de Ultralytics, que la procesa para devolver las detecciones. La aplicación no
guarda esas fotografías ni las poses de ARCore en Firestore, y la imagen
temporal del caché se elimina al reemplazarla o al cerrar la vista.

\subsubsection{Descripción de la interfaz de usuario}

La presente figura~\ref{fig:navegacion} muestra el recorrido del usuario desde la
pantalla de bienvenida hasta las secciones principales de la aplicación.
Después de iniciar sesión, el usuario accede a Inicio si su perfil está
completo. En caso contrario, primero completa la configuración del perfil.
La barra de navegación permite acceder a cuatro pestañas: Inicio, Máquinas,
Progreso y Nutrición. Desde Inicio se accede a las rutinas y al perfil.
Máquinas reúne el escáner y las fichas de consulta, Progreso presenta el
historial y los reportes, y Nutrición contiene el diario de alimentación.

\begin{figure}[H]
\centering
\begin{tikzpicture}[
    font=\footnotesize,
    nvCaja/.style={draw=blue!65!black, fill=blue!8, rounded corners=3pt,
        align=center, inner sep=3pt, minimum height=9mm},
    nvPaso/.style={nvCaja, fill=green!8, text width=3.0cm, font=\scriptsize,
        minimum height=14mm},
    nvPestana/.style={nvCaja, fill=blue!20, text width=2.5cm,
        font=\footnotesize\bfseries},
    nvPantalla/.style={nvCaja, fill=white, text width=2.5cm, font=\scriptsize,
        minimum height=17mm},
    nvLinea/.style={draw=blue!65!black, semithick},
    nvFlujo/.style={draw=blue!65!black, semithick, -{Stealth[length=2mm]}},
    nvAux/.style={draw=blue!65!black, semithick, dashed,
        -{Stealth[length=2mm]}},
    nvGrupo/.style={draw=black!45, dashed, rounded corners=5pt, inner sep=6pt}
]

% ---------- Acceso (sin sesión) ----------
\node[nvCaja, text width=2.0cm] (bienv) at (-6.9,0)
    {\textbf{Bienvenida}\\{\scriptsize Botón Empezar}};
\node[nvCaja, text width=3.9cm] (acceso) at (-3.3,0)
    {\textbf{Acceso}\\{\scriptsize Inicio de sesión o registro (correo o Google)}};
\node[nvCaja, fill=yellow!15, text width=4.6cm] (verif) at (2.3,0)
    {\textbf{Verificación del perfil}\\
     {\scriptsize Guardas de ruta (\texttt{Stack.Protected})}};
\draw[nvFlujo] (bienv) -- (acceso);
\draw[nvFlujo] (acceso) -- (verif);

% ---------- Onboarding (perfil incompleto) ----------
\node[nvPaso] (ob1) at (-5.4,-3.0)
    {\textbf{1. Objetivo y nivel}\\Objetivo y nivel de experiencia};
\node[nvPaso] (ob2) at (-1.8,-3.0)
    {\textbf{2. Datos físicos}\\Edad, altura, peso y días disponibles};
\node[nvPaso] (ob3) at (1.8,-3.0)
    {\textbf{3. Primera rutina}\\Generar con IA, crear personalizada u omitir};
\draw[nvFlujo] (ob1) -- (ob2);
\draw[nvFlujo] (ob2) -- (ob3);
\node[nvGrupo, fit=(ob1)(ob2)(ob3),
    label={[font=\scriptsize]below:Onboarding (perfil incompleto)}] {};
\draw[nvFlujo] (verif.south) -- ++(0,-0.65) -|
    node[pos=0.2, above, font=\scriptsize]{Perfil incompleto} (ob1.north);

% ---------- Aplicación (perfil completo) ----------
\node[font=\scriptsize\bfseries, anchor=north west] (appTitulo)
    at (-7.8,-4.95) {Aplicación (perfil completo)};
\node[nvCaja, fill=gray!12, text width=4.5cm] (perfil) at (-5.3,-6.4)
    {\textbf{Perfil}\\{\scriptsize Editar objetivo, nivel, género, edad y
     altura; registrar peso; cerrar sesión}};
\node[nvCaja, text width=4.4cm] (barra) at (1.5,-6.4)
    {\textbf{Barra de navegación inferior}\\{\scriptsize Cinco pestañas}};
\draw[nvFlujo] (ob3.south) -- (ob3.south |- barra.north);
\draw[nvFlujo] (verif.east) --
    node[above, font=\scriptsize]{Perfil completo} (7.6,0) |- (barra.east);
\draw[nvAux] (barra.west) --
    node[above, font=\tiny, align=center]{avatar del\\encabezado} (perfil.east);

% Pestañas
\foreach \i/\n/\x in {1/Inicio/-6.6, 2/Rutinas/-3.3, 3/Máquinas/0,
                      4/Nutrición/3.3, 5/Progreso/6.6}
    \node[nvPestana] (t\i) at (\x,-8.6) {\n};
\coordinate (bus) at ($(barra.south)+(0,-0.55)$);
\draw[nvLinea] (barra.south) -- (bus);
\draw[nvLinea] (t1.north |- bus) -- (t5.north |- bus);
\foreach \i in {1,2,3,4,5}
    \draw[nvFlujo] (t\i.north |- bus) -- (t\i.north);

% Inicio
\node[nvPantalla, below=5mm of t1] (i1)
    {\textbf{Inicio}\\Tarjeta del día, racha semanal y atajo Escanear máquina};
\draw[nvFlujo] (t1) -- (i1);

% Rutinas
\node[nvPantalla, below=5mm of t2] (r1)
    {\textbf{Mis rutinas}\\Rutina activa y otras; al tocar un ejercicio abre su ficha};
\node[nvPantalla, below=4mm of r1] (r2)
    {\textbf{Nueva rutina}\\Generar con IA (vista previa) o crear personalizada};
\draw[nvFlujo] (t2) -- (r1);
\draw[nvFlujo] (r1) -- (r2);

% Máquinas
\node[nvPantalla, below=5mm of t3] (m1)
    {\textbf{Catálogo}\\Búsqueda y filtros por categoría y tipo};
\node[nvPantalla, below=4mm of m1] (m2)
    {\textbf{Ficha de máquina}\\Instrucciones, errores, prevención, GIF, variantes
     y Añadir a mi rutina};
\node[nvPantalla, below=4mm of m2] (m3)
    {\textbf{Escáner}\\Cámara o ARCore; reconocimiento en la nube o local};
\draw[nvFlujo] (t3) -- (m1);
\draw[nvFlujo] (m1) -- (m2);
\draw[nvFlujo] (m3) -- (m2);
\draw[nvAux] (m1.east) -- ++(0.25,0) |- (m3.east);
\draw[nvAux] (m2.west) -- ++(-0.3,0) |- (r1.east);

% Nutrición
\node[nvPantalla, below=5mm of t4] (n1)
    {\textbf{Diario}\\Meta calórica, macronutrientes y alimentos del día};
\node[nvPantalla, below=4mm of n1] (n2)
    {\textbf{Añadir alimento}\\Buscar (USDA), Rápido o Manual};
\node[nvPantalla, below=4mm of n2] (n3)
    {\textbf{Calc. IMC}\\Cálculo y categoría del IMC};
\draw[nvFlujo] (t4) -- (n1);
\draw[nvFlujo] (n1) -- (n2);
\draw[nvAux] (n1.east) -- ++(0.25,0) |- (n3.east);

% Progreso
\node[nvPantalla, below=5mm of t5] (p1)
    {\textbf{Cuerpo}\\Registro de peso, IMC y gráfico del historial};
\node[nvPantalla, below=4mm of p1] (p2)
    {\textbf{Entrenos}\\Racha, días completados y rutina activa};
\node[nvPantalla, below=4mm of p2] (p3)
    {\textbf{Reporte en PDF}\\Peso, IMC, entrenos y alimentación};
\draw[nvFlujo] (t5) -- (p1);
\draw[nvFlujo] (p1) -- (p2);
\draw[nvFlujo] (p2) -- (p3);

% Entrenar (pantalla completa) y leyenda
\coordinate (base) at ($(m3.south)+(0,-1.0)$);
\node[nvCaja, fill=orange!12, text width=5.6cm] (ent) at ($(base)+(-4.95,0)$)
    {\textbf{Entrenar día} (pantalla completa)\\
     {\scriptsize Series, descansos y registro del entrenamiento}};
\draw[nvFlujo] (i1.south) -- (i1.south |- ent.north);
\draw[nvFlujo] (r2.south) -- (r2.south |- ent.north);
\node[font=\scriptsize, align=left, text width=5.6cm] (leyenda)
    at ($(base)+(4.4,0)$)
    {Línea continua: navegación principal.\\
     Línea discontinua: acción secundaria.};

\node[nvGrupo, inner sep=7pt,
    fit=(appTitulo)(perfil)(barra)(t1)(t5)(i1)(m3)(n3)(p3)(ent)(leyenda)] {};

\end{tikzpicture}
\caption{Esquema de navegación de usuario}
\label{fig:navegacion}
\fuente{Elaboración propia.}
\end{figure}
